%Version 3.1 December 2024
% See section 11 of the User Manual for version history
%
%%%%%%%%%%%%%%%%%%%%%%%%%%%%%%%%%%%%%%%%%%%%%%%%%%%%%%%%%%%%%%%%%%%%%%
%%                                                                 %%
%% Please do not use \input{...} to include other tex files.       %%
%% Submit your LaTeX manuscript as one .tex document.              %%
%%                                                                 %%
%% All additional figures and files should be attached             %%
%% separately and not embedded in the \TeX\ document itself.       %%
%%                                                                 %%
%%%%%%%%%%%%%%%%%%%%%%%%%%%%%%%%%%%%%%%%%%%%%%%%%%%%%%%%%%%%%%%%%%%%%

%%\documentclass[referee,sn-basic]{sn-jnl}% referee option is meant for double line spacing

%%=======================================================%%
%% to print line numbers in the margin use lineno option %%
%%=======================================================%%

%%\documentclass[lineno,pdflatex,sn-basic]{sn-jnl}% Basic Springer Nature Reference Style/Chemistry Reference Style

%%=========================================================================================%%
%% the documentclass is set to pdflatex as default. You can delete it if not appropriate.  %%
%%=========================================================================================%%

%%\documentclass[sn-basic]{sn-jnl}% Basic Springer Nature Reference Style/Chemistry Reference Style

%%Note: the following reference styles support Namedate and Numbered referencing. By default the style follows the most common style. To switch between the options you can add or remove �Numbered� in the optional parenthesis. 
%%The option is available for: sn-basic.bst, sn-chicago.bst%  
 
%%\documentclass[pdflatex,sn-nature]{sn-jnl}% Style for submissions to Nature Portfolio journals
%%\documentclass[pdflatex,sn-basic]{sn-jnl}% Basic Springer Nature Reference Style/Chemistry Reference Style
\documentclass[pdflatex,sn-mathphys-num]{sn-jnl}% Math and Physical Sciences Numbered Reference Style
%%\documentclass[pdflatex,sn-mathphys-ay]{sn-jnl}% Math and Physical Sciences Author Year Reference Style
%%\documentclass[pdflatex,sn-aps]{sn-jnl}% American Physical Society (APS) Reference Style
%%\documentclass[pdflatex,sn-vancouver-num]{sn-jnl}% Vancouver Numbered Reference Style
%%\documentclass[pdflatex,sn-vancouver-ay]{sn-jnl}% Vancouver Author Year Reference Style
%%\documentclass[pdflatex,sn-apa]{sn-jnl}% APA Reference Style
%%\documentclass[pdflatex,sn-chicago]{sn-jnl}% Chicago-based Humanities Reference Style

%%%% Standard Packages
%%<additional latex packages if required can be included here>

\usepackage{graphicx}%
\usepackage{multirow}%
\usepackage{subcaption}
\usepackage{amsmath,amssymb,amsfonts}%
\usepackage{amsthm}%
\usepackage{mathrsfs}%
\usepackage[title]{appendix}%
\usepackage{xcolor}%
\usepackage{textcomp}%
\usepackage{manyfoot}%
\usepackage{booktabs}%
\usepackage{algorithm}%
\usepackage{algorithmicx}%
\usepackage{algpseudocode}%
\usepackage{listings}%
%%%%

%%%%%=============================================================================%%%%
%%%%  Remarks: This template is provided to aid authors with the preparation
%%%%  of original research articles intended for submission to journals published 
%%%%  by Springer Nature. The guidance has been prepared in partnership with 
%%%%  production teams to conform to Springer Nature technical requirements. 
%%%%  Editorial and presentation requirements differ among journal portfolios and 
%%%%  research disciplines. You may find sections in this template are irrelevant 
%%%%  to your work and are empowered to omit any such section if allowed by the 
%%%%  journal you intend to submit to. The submission guidelines and policies 
%%%%  of the journal take precedence. A detailed User Manual is available in the 
%%%%  template package for technical guidance.
%%%%%=============================================================================%%%%

%% as per the requirement new theorem styles can be included as shown below
\theoremstyle{thmstyleone}%
\newtheorem{theorem}{Theorem}%  meant for continuous numbers
%%\newtheorem{theorem}{Theorem}[section]% meant for sectionwise numbers
%% optional argument [theorem] produces theorem numbering sequence instead of independent numbers for Proposition
\newtheorem{proposition}[theorem]{Proposition}% 
%%\newtheorem{proposition}{Proposition}% to get separate numbers for theorem and proposition etc.

\theoremstyle{thmstyletwo}%
\newtheorem{example}{Example}%
\newtheorem{remark}{Remark}%

\theoremstyle{thmstylethree}%
\newtheorem{definition}{Definition}%

\raggedbottom
%%\unnumbered% uncomment this for unnumbered level heads

\begin{document}

\title[Article Title]{An Analysis of Measurement Data and Fire Size Predictions from 15 Years of Christmas Tree Fire Safety Demos}

%%=============================================================%%
%% GivenName	-> \fnm{Joergen W.}
%% Particle	-> \spfx{van der} -> surname prefix
%% FamilyName	-> \sur{Ploeg}
%% Suffix	-> \sfx{IV}
%% \author*[1,2]{\fnm{Joergen W.} \spfx{van der} \sur{Ploeg} 
%%  \sfx{IV}}\email{iauthor@gmail.com}
%%=============================================================%%

\author*[1,2]{\fnm{Isaac T.} \sur{Leventon}}\email{isaac.leventon@nist.gov}

\author[2]{\fnm{Seungjun (Victor)} \sur{Jeong}}\email{sjeong19@terpmail.umd.edu}

\affil*[1]{\orgdiv{Fire Research Division}, \orgname{National Institute of Standards and Technology}, \orgaddress{\street{100 Bureau Drive}, \city{Gaithersburg}, \postcode{20878}, \state{MD}, \country{USA}}}

\affil*[2]{\orgdiv{Department of Fire Protection Engineering}, \orgname{University of Maryland}, \orgaddress{\street{4356 Stadium Dr}, \city{College Park}, \postcode{20742}, \state{MD}, \country{USA}}}

%%==================================%%
%% Sample for unstructured abstract %%
%%==================================%%

\abstract{Each year since 2014, a Christmas tree fire safety event has been organized at the University of Maryland to highlight the potential fire risk of dry Christmas trees. Full-scale fire experiments at these events are used to demonstrate how properly maintaining a cut Christmas tree (to maintain a high moisture content in tree branches and needles) can help to reduce the likelihood of  and to limit fire growth rate and peak fire size. The event is regularly covered by local and national media who help to distribute and amplify the holiday fire safety message. Additionally, the experiments provide in-person experience with large fire dynamics for fire protection engineering students (to support their coursework), they have supported** advanced diagnostics for multiple research projects, and they have enabled an annual, friendly competition that encourages communication and collaboration between members of the fire safety science community. This manuscript introduces a new digital repository that organizes and maintains test data and metadata from 15 years of these experiments. It also quantifies factors impacting fire growth behavior in these tests and offers an analysis of how well competition participants (members of the fire safety science and engineering community) can predict ignition and fire growth behavior of these tree fires.}

%%================================%%
%% Sample for structured abstract %%
%%================================%%

% \abstract{\textbf{Purpose:} The abstract serves both as a general introduction to the topic and as a brief, non-technical summary of the main results and their implications. The abstract must not include subheadings (unless expressly permitted in the journal's Instructions to Authors), equations or citations. As a guide the abstract should not exceed 200 words. Most journals do not set a hard limit however authors are advised to check the author instructions for the journal they are submitting to.
% 
% \textbf{Methods:} The abstract serves both as a general introduction to the topic and as a brief, non-technical summary of the main results and their implications. The abstract must not include subheadings (unless expressly permitted in the journal's Instructions to Authors), equations or citations. As a guide the abstract should not exceed 200 words. Most journals do not set a hard limit however authors are advised to check the author instructions for the journal they are submitting to.
% 
% \textbf{Results:} The abstract serves both as a general introduction to the topic and as a brief, non-technical summary of the main results and their implications. The abstract must not include subheadings (unless expressly permitted in the journal's Instructions to Authors), equations or citations. As a guide the abstract should not exceed 200 words. Most journals do not set a hard limit however authors are advised to check the author instructions for the journal they are submitting to.
% 
% \textbf{Conclusion:} The abstract serves both as a general introduction to the topic and as a brief, non-technical summary of the main results and their implications. The abstract must not include subheadings (unless expressly permitted in the journal's Instructions to Authors), equations or citations. As a guide the abstract should not exceed 200 words. Most journals do not set a hard limit however authors are advised to check the author instructions for the journal they are submitting to.}

\keywords{burning behavior, fire safety, vegetation, wildfire}

%%\pacs[JEL Classification]{D8, H51}

%%\pacs[MSC Classification]{35A01, 65L10, 65L12, 65L20, 65L70}

\maketitle
\clearpage
\section{Introduction}\label{sec:intro}

Every year, fire departments in the United States respond to an average of 150 to 200 Christmas Tree fires~\cite{Campbell2022christmas} (despite their use in a relatively short time each year).Home structure fires involving Christmas trees are, on average, substantially more deadly\footnote{A 2017 the National Fire Protection Association (NFPA) report, \textit{Home Structure Fires Involving Christmas Trees}~\cite{ahrens2017christmastrees}, noted that, "If a Christmas tree fire occurs, the risk of death is unusually high... On average, one of every 32 reported home Christmas tree fires resulted in a death, compared to an average of one death per 143 total reported home fires." This represents a factor of 4.5x increase.} than other house fires~\cite{ahrens2017christmastrees}. Beginning in 2014~\cite{pbsnewshour2014christmastrees}, a series of Christmas Tree fire safety demonstrations has been coordinated each year to: (1) support and promote better public awareness of fire safety during the holiday season, (2) to highlight relevant experimental\cite{Giovanni's work, Gollner/Leventon} and modeling\cite{Mell, Maranghides, McDermott} research projects and tools, and (3) to offer Fire Protection Engineering students the opportunity to experience and measure fire behaviors that they [typically only study on paper during their regular coursework**]. These tests are conducted at the University of Maryland, College Park (UMD) with support from faculty and researchers of the UMD Department of Fire Protection Engineering (FPE) and the National Institute of Standards and Technology (NIST).

In its first year, this fire safety event served as the final for an introductory fire science program taught at UMD (this program was supported in 2015 by a grant from the Society of Fire Protection Engineers, SFPE, Scientific and Educational Foundation). Later, senior- and graduate-level students of the FPE department have watched these fire experiments (from a safe distance) to contextualize the fire behaviors that they study. In various years, these experiments were included in courses such as fire dynamics, enclosure fire modeling, and the graduate fire dynamics laboratory to allow students to practice assessing design fires, calculating fire risk (e.g., zone of influence and heat transfer at a distance), and performing statistical and uncertainty analyses of large fire datasets.

As the risk posed by wildfires and wildland urban interface (WUI) fires continues to grow~\cite{parks2020warmer, cunningham2024increasing,kelley2025state}, test matrices for these fire safety demonstrations have been expanded and additional measurement devices have been incorporated to produce unique datasets (e.g., firebrand yield) and to quantify factors that impact tree burning behavior (e.g., ignition source, ambient conditions, or tree species, size, and moisture content). In some cases, these results have been directly included in conference presentations and/or journal publications~\cite{leventon-gollner}; in others, the experiments provide an opportunity for calibration measurement devices used in other studies~\cite{adetolla-raffan}.

In addition to sharing a fire safety message and supporting research efforts, these experiments also serve as the target for an annual competition to predict the burning behavior (i.e., heat release rate, HRR) of Christmas Tree fires, from ignition to final burnout. This competition provides a platform for open communication between members of the fire safety science community and an opportunity to encourage and motivate current researchers and the next generation of scientists and engineers to explore and develop collaborations related to emerging areas of fire safety science. 

This paper highlights the development of these fire safety demos from a small, local event within the Department of Fire Protection Engineering at UMD to annual tradition with a global footprint: each year, the Christmas Tree Fire Size Competition includes hundreds predictions from participants in dozens of countries across the world and [related fire safety multi-media/videos] have received millions of impressions while the fire safety experiments receive continued local and national media attention. Specifically, this manuscript describes the development, organization, and standard formatting used in a new online repository 
of the measurement data and test metadata obtained during these experiments. Corresponding analysis of these results highlights impacts of key factors that affect observed burning behavior. A review of competition results similarly highlights how the HRR prediction competition is organized, how it has grown, and what can be learned from these predictions (i.e., what are participants - all members of the fire safety science and engineering community - good or bad at predicting and (how) has this evolved?).


\section{Experimental}\label{sec:exp}
\subsection{Background}
Christmas tree fire safety research has been conducted for decades, demonstrating the impact of moisture content~\cite{stroup1999scotch,damant1994christmas,Babrauskas2001, white1997flammability}, geometry~\cite{madrzykowski2008impact}, tree size~\cite{evans2004physics}, and even the impact of residential sprinklers\cite{madrzykowski2008impact} on tree burning behavior. More recently, different physics-based models (WFDS~\cite{Mell2009}, FireStar3D~\cite{Accary2024}, FDS 6.9.1~\cite{Gupta2025}) have been validated to predict experimentally-measured burning behavior of 1.5~m to 5~m tall trees, with further analysis exploring model sensitivity: to grid resolution, definition of ignition conditions, fuel moisture content, representation of fuel particles (geometry, particle size distribution, heats of combustion), and description of pyrolysis/oxidation reactions. Related experimental studies of the burning behavior of large trees and other vegetative fuels (whole or components)~\cite{Gio-3x, catchpole1998rate, Babrauskas2001, dahale2013effects, Morandini2019, etlinger2004development, white1997flammability} further highlight the coupled influence of tree moisture content, ignition source (size, strength, and location), ambient wind, and tree size, geometry, and bulk density on ignitability, early fire growth rate, and peak fire size. 
%Morandini2019 - beautiful experimental characterization of particle size / mass distribution as a function of shrub height. exp/modeling paper. vary ignition location - lower in shrub --> larger fires

In 2020, an experiments were conducted at NIST to examine the difference in fire hazard between a dry and a well-watered Christmas tree  ~\cite{hoehler2020fire}, with the explicit goal of recording both accurate HRR measurements and high-quality video recordings (in a realistic setting) for fire safety awareness. 

Fire safety messaging: let's get a review on fear, fear + action.\\
don't forget that ACS publication on tree fire safety. Highlight value/risk in repeated information (borax as fire retardant..)

From J Bonny - some background/refs on what experience(s) help best build awareness of potential fire safety risk.


Numerous works reviews have highlights xyz that are important gaps or strong correlation/impact on key burning behavior.\\

Repeat the value of this exercise is not simply the dataset, but the unique ways that it can be shared, ingested, worked with, etc..




%"Recent experiments involving over 20 conifer species have been conducted at Washington State University to determine the MC where branches initially begin to fail an ignition test, and the MC for consistent failure (Figure 3). During these studies, 15.3 cm-long branch samples were exposed to a flame from a small alcohol lamp for 5 seconds. If the sample failed to burn or selfextinguished without any additional spread when the sample was removed from the flame, it passed the test. If there was any spread of the flame after the sample was removed from the flame, it failed the test. Results from these tests indicates that there is considerable variation in the MC where different species failed this test. In addition, there is often a large transition zone of MC from the point of initially failing the flammability test to the point where samples consistently fail. There is also a difference in the MC where samples fail the test and the MC when samples are totally consumed by the fire. For example, Douglas-fir begins to fail the branch flammability test at about 68\% MC, but is totally consumed only when dried to about 30\% MC (Figure 4)."[]~\cite{Babrauskas2001} Chastanger's section]

%Impact of ignition source:
%(on PVC trees) "The specimens proved impossible to ignite with a small flame, but using the UL procedure (crumpled newspaper in a 61cm circle around base of tree) they produced rapidly developing fires, with 500 kW being attained 11?20 s after ignition..." substantially larger fire when paper ignited in four locations [This come from unpublished work included in Vyto's Heat Release Rates chapter in the SFPE handbook]
%Ref\cite{Accary2024} explored how burner power and duration affected MLR predictions (only )

%%%%%%%%%%%
\subsection{Measurements and Test Procedure}\label{ssec:exp-measurements}
After procurement, the initial mass, height, and weight of each tree is measured and then trees are kept unwatered for approximately 14 to 21 days in laboratory storage at the University of Maryland\footnote{Each year, at least one tree is kept in a 5 gallon bucket of water (continuously refilled) to allow for comparison of well-watered and dry tree burning behavior. In some years, additional tests are conducted with expanded conditioning durations to explore a wider range of tree moisture contents. All trees are stored approximately at approximately 40\% relative humidity and 20$^\circ$C}. Immediately prior to testing, small branches are removed from the base of each tree (approximately 100~g to 200~g), weighed, dried in an oven at 80$^\circ$C for approximately 2 days, and then re-weighed. Tree moisture content is calculated as $MC=\frac{m_{0}-m_{\text{dry}}}{m_\text{{dry}}}$.%, where the subscripts '0' and 'dry' indicate branch masses before and after oven drying.
The 'dry' mass of the full tree is then measured, the tree is secured in a metal stand, and then this stand is placed above a mass balance that is protected by a radiation shield and several layers of rigid thermal insulation. 

During tests, the mass balance/tree are positioned at the center of (underneath) an exhaust hood that measures 3.66~m by 3.66~m (square profile, sides nominally 12 ft wide). The hood is rated for a nominal air flow rate of 10,000~CFM. Three sides of the exhaust hood were enclosed using fire-rated fabric curtains to help contain the potential spread of firebrands generated during testing.  When used, heat flux gauges and/or thermocouples are then secured in place and their locations are measured/recorded. The ignition source is then moved into position below the lowest branches of the tree, data and video recording are started and confirmed, and the ignition source is ignited. Once the lower branches of the tree ignite, the exhaust hood is turned on. Trees are allowed to burn until self-extinction, after which final tree mass is measured, data and video recording were stopped, and pictures of the remaining sample were taken.

\subsubsection{Ignition Source}
Previous studies have demonstrated how ignition source type, strength, duration, and location all can affect fire growth characteristics of both vegetative fuels and materials/products found in the built environment~\cite{dembsey1993effect,ai1987effects,cleary1994influence, Babrauskas2001, Morandini2019, Accary2024, gould2022initial, Babrauskas1985Fire}. In a wildfire context, multiple studies have also assessed the impact of ignition source (e.g., point source vs. smaller or larger line fires) on wildfire spread rate~\cite{Rothermel1984Aerial, Johansen1987Ignition, Alexander1991PorterLake,gould2022initial}. 
Though this has been explored in some modeling studies (e.g., Ref~\cite{Accary2024}, limited information is available from experiments designed to systematically quantify how variations in ignition size/severity, duration, and/or location affect fire growth rate/peak fire size of individual vegetative fuels (which in turn controls/can strongly impact ROS).

With this in mind, careful consideration was given to the ignition sources used in this Christmas Tree fire safety study. This ensured that ignition could be performed repeatably and that the source reasonably represented a small fire hazard found in the home. Most often, this meant filling a 6~cm to 10~cm disc (cut from rigid ceramic insulation board and backed/wrapped by aluminum foil) with approximately 30~mL of heptane and placing this beneath the lowest branches at the base of the tree. Both to conduct parametric studies and to make the fire size prediction competition more challenging (i.e., such that ignition and fire growth could be reasonably expected to change each year) the location of this small pool fire could be moved from year to year (e.g., near the center of the tree or towards the outer branches). Also, in some years, this pool fire was replaced by a full ring (cut from the same ceramic insulation) to allow for uniform ignition around the entire base of the tree. Additionally, in later years of the competition, tree ignition was accomplished by the spread of a ground fire across an angled bed of excelsior placed below the base of the tree - this new ignition source both (1) offered a new connection to wildland fire scenarios and (2) emphasized prediction of tree ignition time delay (due to flame spread) as an important scoring category.

%%%%Impact of ignition source on fire growth rate:
%" effect of ignition source exposure and specimen configuration on the actual fire growth characteristics of polyvinyl chloride (PVC) foam wall covering... :(i) the PVC foam has a ?critical? ignition source strength of 64?75 kW; (ii) preheating the PVC foam greatly increases its peak net rate of heat release."~\cite{dembsey1993effect}

%"[really not the best - different burner sizes, HRR in a room lined with chipboard (hink MDF). burner HRR controls time to flashover / T=500C]~\cite{ai1987effects}

%" upholstered chairs ... five different fabric/foam combinations ...ignited by a cigarette, a small match-like flame, an incandescent lamp, a space heater, and a large flame source (TB 133 ignition source)... for any chair type (more/less variation in t\_peak for diff chairs), the time to the peak heat release rate depended on the ignition sequence, but the magnitude of the peak did not"~\cite{cleary1994influence}


%"Fifty-eight fires of three ignition patterns (point, 400-mm line and 800-mm line) were carried out at two air speeds (1.25 and 2.0?m s?1) and two dead fuel moisture content (FMC) groups for each air speed (?7.5~\% and >7.5\% oven-dry weight for the low air speed and ?5\% and >5\% for the high air speed)."[shows how time to steady/peak spread rate of litter fuel beds is affected by ignition source \cite{gould2022initial}

%"A NIST study examined various ignition sources, ranging from 5 W to over 300 kW [45]. The sources included both realistic igniting objects (cigarettes, matches, burning paper lunch bags, etc.) and schematic ones (small gas burners and wood cribs). It was found that, as the power output of the ignition source increased, the peak heat flux generally did not increase. Instead, only the area covered by the peak heat flux progressively increased" ... "This study also looks at British, NZ, American standards for igniting upholstered furniture using different sources"
\subsubsection{Measurement Devices}
\textbf{Burning Rate and Heat Release}\\
Measurements recorded during this ongoing test series include: initial and final tree masses, time-resolved mass loss rate during burning, and heat flux at a distance. In certain test series, local temperature measurements (in or near the tree during burning) or total firebrand yield and mass/projected areas of individual firebrands were also recorded. In general, all experiments were conducted in the large fire lab at the University of Maryland. However, in 2024, the competition was based on measurements of extremely large (5.8~m** tall) tree fires that were performed in the National Fire Research Laborartory (NFRL) at NIST. This measurements and corresponding uncertainties of this dataset -- including MLR**, HRR, heat flux at \#** locations -- is detailed elsewhere~\cite{Gio et al.}. Excluding this dataset, all measurements were recorded as described below.

A Mettler Toldeo **Model Number* mass balance \footnote{ The identification of any commercial product or trade name does not imply endorsement or recommendation by NIST (or any other contributing institution)} with a 32~kg range and 0.1~g precision was used for all\footnote{In Winter 2019/2020, a detailed study of firebrand generation was conducted; the masses of individual firebrands was measured using a balance with higher precision, as described below.} mass measurements obtained in this test series (both static, and time-resolved). Sample mass was continuously measured at 1~Hz throughout the duration of each experiment. Mass loss rate was then calculated at each time step, $t_i$, as the numerical derivative of sample mass curves across a $\pm$~2~s interval, as follows: 
\begin{equation}
	\dot{m}_{i} = \frac{\left(m_{i-2}-m_{i+2}\right)}{\left(t_{i+2}-t_{i-2}\right)}
\end{equation}

For the purposes of the competition, heat release rate was calculated as the product of measured mass loss rate and an assumed heat of combustion, typically $\Delta H_c=17.7$~kJ/g. 

Heat Flux: ()approximately 2~m away from the outside edge of the tree


\textbf{Firebrand Generation}\\
"Surrounding each tree, a series of
approximately 60 water-filled aluminum trays covered the floor space inside this enclosure
to extinguish and collect firebrands formed during testing.
Inside the exhaust hood, approximately 3 m above the floor, a 2.5 m by 2.5 m square
frame was built to support a steel mesh (opening dimensions, 0.25 cm square). This mesh
was installed to contain any firebrands that were formed and lofted upwards during burning.
The total open area of this mesh screen was sufficiently large such that air flow rate through
the mesh was low enough that even small leaves (high surface area to volume/mass ratio)
would not stick to the mesh due to operation of the exhaust hood. However, while trees were
burning, due to the increased fire-induced upwards flow, some firebrands were captured by
this screen. Typically, these firebrands were small, lightweight leaves/needles that would
smolder to ash within 45 s (i.e., by the end of the experiment).


All firebrands generated during testing were carefully collected by filtering each waterfilled pan, draining the water, and saving the remaining pieces in a new pan for further
sorting. Firebrands were sorted into two groups ? smaller brands (approximately 10 cm
or smaller), which were charred across the majority of their surface (and thus assumed
to be smoldering prior to collection in the water pans) and larger pieces (e.g., branches)
that had broken off near their base (at attachment point to the trunk) but remained mostly
unburned. All firebrands were dried for 48 hours in a convection oven at approximately
120$^\circ$C and then weighed; measured masses of smaller and branch-type firebrands were
noted as mfirebrands and mbranch, respectively.


In total, more than 10,700 individual firebrands were collected in this series of experiments; after drying, the mass and projected area of each one was measured as follows.
First, individual firebrands were selected by ceramic tweezers and placed on a Sartorius
Cubis MSE3203S mass balance (3200 g range; 1 mg repeatability, i.e., standard deviation
when calculated from 5 cycles under the conditions reported in Ref. [39]). After weighing,
firebrands were placed (12 at a time) onto predetermined regions of a 0.64 m by 0.73 m
sheet of white poster board. A photograph of these firebrands was then taken using a Nikon
D3200 DSLR camera with a Nikon AF-S DX NIKKOR 18-55mm lens that was positioned
approximately 0.8 m above the board."


\subsection{Materials/Fuels and Sample Conditioning}\label{ssec:exp-matls}
As of 2026, \#\# individual trees of \# species of trees have been tested as part of this ongoing Christmas Tree fire test series, each with a range of initial moisture contents. These tree species were selected because of their widespread presence across different regions in North America and the relative ease with which they could be acquired around the holiday season. Each year, typically in late November or early December, trees were purchased (and in some cases, received as donations) from hardware supply stores and/or tree farms in Central Maryland. Figure~\ref{fig:3trees} shows representative images of three of these tree species prior to testing (each studied during the winter of 2019/2020~\cite{leventon-gollner}). The common and scientific names of each of all tree species tested to date, and their respective habitats, are identified in Table~\ref{tab:tree-species}. 

\begin{table}[]
	\footnotesize
	\caption{Species Tested}
	\label{tab:tree-species}
	\begin{tabular}{p{2.5cm}|p{2.5cm}p{2.5cm}p{2.5cm}p{2.5cm}}
		\hline
		\textbf{Common Name}&Balsam Fir&Canaan Fir &Douglas-Fir&Fraser Fir\\
		\textbf{Scientific Name}&Abies Balsamea**&???**&Pseudotsuga Menziesii&Abies Fraseri\\
		\textbf{Habitat}&Eastern and Central Canada, Northeastern US &&Rocky Mountain Region, North America&Southeastern Appalachian Mountains \\
		\hline
		\textbf{Height}&&&&\\
		\textbf{Weight}&&&&\\
		\textbf{Moisture Content}&&&&\\
		\textbf{Ignition Source}&&&&\\
		%\multicolumn{3}{p{14.4cm}}{\textsuperscript{A} Rocky Mountain Research Station (Missoula, Montana)}\\
%		\multicolumn{5}{p{14.4cm}}{\textsuperscript{B} Peat collected near Prudhoe bay was obtained from a flooded wetland meadow (5~cm to 10~cm standing water depth) dominated by eriophorum species.}\\
	\end{tabular}
\end{table}

\begin{figure} 
	\centering 
	\begin{subfigure}[t]{0.3\textwidth}
		\centering
		\includegraphics[width=\textwidth]{../FIGURES/BalsamFir_inlab.png}
		\caption{Balsam fir, \textit{Abies Balsamea}; Eastern and Central Canada, Northeastern US }
		\label{Fig:BalsamFir}
	\end{subfigure}
	\begin{subfigure}[t]{0.3\textwidth}
		\centering
		\includegraphics[width=\textwidth]{../FIGURES/Douglas-Fir_inlab.png}
		\caption{Douglas-Fir, \textit{Pseudotsuga Menziesii}; Rocky Mountain Region, North America}
		\label{Fig:DouglasFir}
	\end{subfigure}
	\begin{subfigure}[t]{0.3\textwidth}
		\centering
		\includegraphics[width=\textwidth]{../FIGURES/FraserFir_inlab.png}
		\caption{Fraser Fir, \textit{Abies Fraseri}; Southeastern Appalachian Mountains }
		\label{Fig:FraserFir}
	\end{subfigure}
	\caption{Representative images of tree species prior to burning. }
	\label{fig:3trees}
\end{figure}


\section{HRR Prediction Competition}\label{sec:competition}
\subsection{Development and Organization}\label{ssec:comp-organization}
General organization/timeline (practical, media, etc.). What are we asked to predict, why, scoring metrics, prizes
\subsection{Preparing, Analyzing, and Scoring Predictions}\label{ssec:comp-predictions}
How do we generate these predictions/how has it evolved? Growth in participation / acknowledgment of support

\section{Results and Discussion}\label{sec:results}
\subsection{Experimental Measurements}\label{ssec:results-experimental}
\subsection{Analysis of Competition Predictions}\label{ssec:results-competition}
\subsection{Fire Safety Messaging/multimedia}\label{ssec:results-safety}

Discussions should be brief and focused. In some disciplines use of Discussion or `Conclusion' is interchangeable. It is not mandatory to use both. Some journals prefer a section `Results and Discussion' followed by a section `Conclusion'. Please refer to Journal-level guidance for any specific requirements. 

\section{Conclusion}\label{sec13}

In the 15 years since, these demos have been repeated allowing for 
\begin{itemize}
	\item Students can apply xyz from class/real world experience of what numbers/terms look/feel like
	\item Friendly competition to foster ABC (from 15x submission all at UMD to 100s per year across \#\#  states in total, \#\#  countries in total)
	\item Fire safety demos (yearly for local media, prepared packages now shown in all states in the US)
	\item Proper measurement science studies that explore the impact of MC, ignition source/size/location, wind on ignitabiilty, fire growth rate, peak fire size, and firebrand production/yield
\end{itemize}

We burned a lot of trees. Moisture content is critical to ignitability/fire growth. Ignition/air flow has a strong impact (2x) on peak fire size/growth rate.\\
As a community we're surprisingly good at predicting X and Y but struggle with 1, 2, and 3.



\section{Equations}\label{sec4}

Equations in \LaTeX\ can either be inline or on-a-line by itself (``display equations''). For
inline equations use the \verb+$...$+ commands. E.g.: The equation
$H\psi = E \psi$ is written via the command \verb+$H \psi = E \psi$+.

For display equations (with auto generated equation numbers)
one can use the equation or align environments:
\begin{equation}
\|\tilde{X}(k)\|^2 \leq\frac{\sum\limits_{i=1}^{p}\left\|\tilde{Y}_i(k)\right\|^2+\sum\limits_{j=1}^{q}\left\|\tilde{Z}_j(k)\right\|^2 }{p+q}.\label{eq1}
\end{equation}
where,
\begin{align}
D_\mu &=  \partial_\mu - ig \frac{\lambda^a}{2} A^a_\mu \nonumber \\
F^a_{\mu\nu} &= \partial_\mu A^a_\nu - \partial_\nu A^a_\mu + g f^{abc} A^b_\mu A^a_\nu \label{eq2}
\end{align}
Notice the use of \verb+\nonumber+ in the align environment at the end
of each line, except the last, so as not to produce equation numbers on
lines where no equation numbers are required. The \verb+\label{}+ command
should only be used at the last line of an align environment where
\verb+\nonumber+ is not used.
\begin{equation}
Y_\infty = \left( \frac{m}{\textrm{GeV}} \right)^{-3}
    \left[ 1 + \frac{3 \ln(m/\textrm{GeV})}{15}
    + \frac{\ln(c_2/5)}{15} \right]
\end{equation}
The class file also supports the use of \verb+\mathbb{}+, \verb+\mathscr{}+ and
\verb+\mathcal{}+ commands. As such \verb+\mathbb{R}+, \verb+\mathscr{R}+
and \verb+\mathcal{R}+ produces $\mathbb{R}$, $\mathscr{R}$ and $\mathcal{R}$
respectively (refer Subsubsection~\ref{subsubsec2}).

\section{Tables}\label{sec5}

Tables can be inserted via the normal table and tabular environment. To put
footnotes inside tables you should use \verb+\footnotetext[]{...}+ tag.
The footnote appears just below the table itself (refer Tables~\ref{tab1} and \ref{tab2}). 
For the corresponding footnotemark use \verb+\footnotemark[...]+

\begin{table}[h]
\caption{Caption text}\label{tab1}%
\begin{tabular}{@{}llll@{}}
\toprule
Column 1 & Column 2  & Column 3 & Column 4\\
\midrule
row 1    & data 1   & data 2  & data 3  \\
row 2    & data 4   & data 5\footnotemark[1]  & data 6  \\
row 3    & data 7   & data 8  & data 9\footnotemark[2]  \\
\botrule
\end{tabular}
\footnotetext{Source: This is an example of table footnote. This is an example of table footnote.}
\footnotetext[1]{Example for a first table footnote. This is an example of table footnote.}
\footnotetext[2]{Example for a second table footnote. This is an example of table footnote.}
\end{table}

\noindent
The input format for the above table is as follows:

%%=============================================%%
%% For presentation purpose, we have included  %%
%% \bigskip command. Please ignore this.       %%
%%=============================================%%
\bigskip
\begin{verbatim}
\begin{table}[<placement-specifier>]
\caption{<table-caption>}\label{<table-label>}%
\begin{tabular}{@{}llll@{}}
\toprule
Column 1 & Column 2 & Column 3 & Column 4\\
\midrule
row 1 & data 1 & data 2	 & data 3 \\
row 2 & data 4 & data 5\footnotemark[1] & data 6 \\
row 3 & data 7 & data 8	 & data 9\footnotemark[2]\\
\botrule
\end{tabular}
\footnotetext{Source: This is an example of table footnote. 
This is an example of table footnote.}
\footnotetext[1]{Example for a first table footnote.
This is an example of table footnote.}
\footnotetext[2]{Example for a second table footnote. 
This is an example of table footnote.}
\end{table}
\end{verbatim}
\bigskip
%%=============================================%%
%% For presentation purpose, we have included  %%
%% \bigskip command. Please ignore this.       %%
%%=============================================%%

\begin{table}[h]
\caption{Example of a lengthy table which is set to full textwidth}\label{tab2}
\begin{tabular*}{\textwidth}{@{\extracolsep\fill}lcccccc}
\toprule%
& \multicolumn{3}{@{}c@{}}{Element 1\footnotemark[1]} & \multicolumn{3}{@{}c@{}}{Element 2\footnotemark[2]} \\\cmidrule{2-4}\cmidrule{5-7}%
Project & Energy & $\sigma_{calc}$ & $\sigma_{expt}$ & Energy & $\sigma_{calc}$ & $\sigma_{expt}$ \\
\midrule
Element 3  & 990 A & 1168 & $1547\pm12$ & 780 A & 1166 & $1239\pm100$\\
Element 4  & 500 A & 961  & $922\pm10$  & 900 A & 1268 & $1092\pm40$\\
\botrule
\end{tabular*}
\footnotetext{Note: This is an example of table footnote. This is an example of table footnote this is an example of table footnote this is an example of~table footnote this is an example of table footnote.}
\footnotetext[1]{Example for a first table footnote.}
\footnotetext[2]{Example for a second table footnote.}
\end{table}

In case of double column layout, tables which do not fit in single column width should be set to full text width. For this, you need to use \verb+\begin{table*}+ \verb+...+ \verb+\end{table*}+ instead of \verb+\begin{table}+ \verb+...+ \verb+\end{table}+ environment. Lengthy tables which do not fit in textwidth should be set as rotated table. For this, you need to use \verb+\begin{sidewaystable}+ \verb+...+ \verb+\end{sidewaystable}+ instead of \verb+\begin{table*}+ \verb+...+ \verb+\end{table*}+ environment. This environment puts tables rotated to single column width. For tables rotated to double column width, use \verb+\begin{sidewaystable*}+ \verb+...+ \verb+\end{sidewaystable*}+.

\begin{sidewaystable}
\caption{Tables which are too long to fit, should be written using the ``sidewaystable'' environment as shown here}\label{tab3}
\begin{tabular*}{\textheight}{@{\extracolsep\fill}lcccccc}
\toprule%
& \multicolumn{3}{@{}c@{}}{Element 1\footnotemark[1]}& \multicolumn{3}{@{}c@{}}{Element\footnotemark[2]} \\\cmidrule{2-4}\cmidrule{5-7}%
Projectile & Energy	& $\sigma_{calc}$ & $\sigma_{expt}$ & Energy & $\sigma_{calc}$ & $\sigma_{expt}$ \\
\midrule
Element 3 & 990 A & 1168 & $1547\pm12$ & 780 A & 1166 & $1239\pm100$ \\
Element 4 & 500 A & 961  & $922\pm10$  & 900 A & 1268 & $1092\pm40$ \\
Element 5 & 990 A & 1168 & $1547\pm12$ & 780 A & 1166 & $1239\pm100$ \\
Element 6 & 500 A & 961  & $922\pm10$  & 900 A & 1268 & $1092\pm40$ \\
\botrule
\end{tabular*}
\footnotetext{Note: This is an example of table footnote this is an example of table footnote this is an example of table footnote this is an example of~table footnote this is an example of table footnote.}
\footnotetext[1]{This is an example of table footnote.}
\end{sidewaystable}

\section{Figures}\label{sec6}

As per the \LaTeX\ standards you need to use eps images for \LaTeX\ compilation and \verb+pdf/jpg/png+ images for \verb+PDFLaTeX+ compilation. This is one of the major difference between \LaTeX\ and \verb+PDFLaTeX+. Each image should be from a single input .eps/vector image file. Avoid using subfigures. The command for inserting images for \LaTeX\ and \verb+PDFLaTeX+ can be generalized. The package used to insert images in \verb+LaTeX/PDFLaTeX+ is the graphicx package. Figures can be inserted via the normal figure environment as shown in the below example:

%%=============================================%%
%% For presentation purpose, we have included  %%
%% \bigskip command. Please ignore this.       %%
%%=============================================%%
\bigskip
\begin{verbatim}
\begin{figure}[<placement-specifier>]
\centering
\includegraphics{<eps-file>}
\caption{<figure-caption>}\label{<figure-label>}
\end{figure}
\end{verbatim}
\bigskip
%%=============================================%%
%% For presentation purpose, we have included  %%
%% \bigskip command. Please ignore this.       %%
%%=============================================%%

\begin{figure}[h]
\centering
\includegraphics[width=0.9\textwidth]{fig.eps}
\caption{This is a widefig. This is an example of long caption this is an example of long caption  this is an example of long caption this is an example of long caption}\label{fig1}
\end{figure}

In case of double column layout, the above format puts figure captions/images to single column width. To get spanned images, we need to provide \verb+\begin{figure*}+ \verb+...+ \verb+\end{figure*}+.

For sample purpose, we have included the width of images in the optional argument of \verb+\includegraphics+ tag. Please ignore this. 

\section{Cross referencing}\label{sec8}

Environments such as figure, table, equation and align can have a label
declared via the \verb+\label{#label}+ command. For figures and table
environments use the \verb+\label{}+ command inside or just
below the \verb+\caption{}+ command. You can then use the
\verb+\ref{#label}+ command to cross-reference them. As an example, consider
the label declared for Figure~\ref{fig1} which is
\verb+\label{fig1}+. To cross-reference it, use the command 
\verb+Figure \ref{fig1}+, for which it comes up as
``Figure~\ref{fig1}''. 

To reference line numbers in an algorithm, consider the label declared for the line number 2 of Algorithm~\ref{algo1} is \verb+\label{algln2}+. To cross-reference it, use the command \verb+\ref{algln2}+ for which it comes up as line~\ref{algln2} of Algorithm~\ref{algo1}.

\subsection{Details on reference citations}\label{subsec7}

Standard \LaTeX\ permits only numerical citations. To support both numerical and author-year citations this template uses \verb+natbib+ \LaTeX\ package. For style guidance please refer to the template user manual.

Here is an example for \verb+\cite{...}+: \cite{bib1}. Another example for \verb+\citep{...}+: \citep{bib2}. For author-year citation mode, \verb+\cite{...}+ prints Jones et al. (1990) and \verb+\citep{...}+ prints (Jones et al., 1990).

All cited bib entries are printed at the end of this article: \cite{bib3}, \cite{bib4}, \cite{bib5}, \cite{bib6}, \cite{bib7}, \cite{bib8}, \cite{bib9}, \cite{bib10}, \cite{bib11}, \cite{bib12} and \cite{bib13}.


\section{Methods}\label{sec11}

Topical subheadings are allowed. Authors must ensure that their Methods section includes adequate experimental and characterization data necessary for others in the field to reproduce their work. Authors are encouraged to include RIIDs where appropriate. 

\textbf{Ethical approval declarations} (only required where applicable) Any article reporting experiment/s carried out on (i)~live vertebrate (or higher invertebrates), (ii)~humans or (iii)~human samples must include an unambiguous statement within the methods section that meets the following requirements: 

\begin{enumerate}[1.]
\item Approval: a statement which confirms that all experimental protocols were approved by a named institutional and/or licensing committee. Please identify the approving body in the methods section

\item Accordance: a statement explicitly saying that the methods were carried out in accordance with the relevant guidelines and regulations

\item Informed consent (for experiments involving humans or human tissue samples): include a statement confirming that informed consent was obtained from all participants and/or their legal guardian/s
\end{enumerate}

If your manuscript includes potentially identifying patient/participant information, or if it describes human transplantation research, or if it reports results of a clinical trial then  additional information will be required. Please visit (\url{https://www.nature.com/nature-research/editorial-policies}) for Nature Portfolio journals, (\url{https://www.springer.com/gp/authors-editors/journal-author/journal-author-helpdesk/publishing-ethics/14214}) for Springer Nature journals, or (\url{https://www.biomedcentral.com/getpublished/editorial-policies\#ethics+and+consent}) for BMC.

\backmatter

\bmhead{Supplementary information}

If your article has accompanying supplementary file/s please state so here. 

\bmhead{Acknowledgements}

These experiments have been made possible thanks to the efforts of researchers, technicians, support staff, and students at both the University of Maryland (UMD) and the National Institute of Standards and Technology (NIST). Their continuous support has enabled yearly testing, recording of measurement data, competition coordination, and media outreach / fire safety messaging. In particular, the following student researchers are thanked for their support in data archiving, formatting, and analysis:  Seungjun (Victor) Jeong, Mahdi Tlemsani**. A special thanks is also made to the authors of a recent large-scale fire experiments conducted at NIST on Douglas-Fir trees between 2.1~m and 5.8~m tall; use of their mass loss measurement data and thoughtful discussions improved the analysis included in this manuscript.


\section*{Declarations}

All measurement data presented in this manuscript is found online in a freely-available, version-controlled github repository: \url{https://github.com/leventon/Christmas_Tree_Fires}


%%===================================================%%
%% For presentation purpose, we have included        %%
%% \bigskip command. Please ignore this.             %%
%%===================================================%%

\begin{appendices}

\section{Section title of first appendix}\label{secA1}

An appendix contains supplementary information that is not an essential part of the text itself but which may be helpful in providing a more comprehensive understanding of the research problem or it is information that is too cumbersome to be included in the body of the paper.

%%=============================================%%
%% For submissions to Nature Portfolio Journals %%
%% please use the heading ``Extended Data''.   %%
%%=============================================%%

%%=============================================================%%
%% Sample for another appendix section			       %%
%%=============================================================%%

%% \section{Example of another appendix section}\label{secA2}%
%% Appendices may be used for helpful, supporting or essential material that would otherwise 
%% clutter, break up or be distracting to the text. Appendices can consist of sections, figures, 
%% tables and equations etc.

\end{appendices}

%%===========================================================================================%%
%% If you are submitting to one of the Nature Portfolio journals, using the eJP submission   %%
%% system, please include the references within the manuscript file itself. You may do this  %%
%% by copying the reference list from your .bbl file, paste it into the main manuscript .tex %%
%% file, and delete the associated \verb+\bibliography+ commands.                            %%
%%===========================================================================================%%

\bibliography{References.bib}% common bib file
%% if required, the content of .bbl file can be included here once bbl is generated
%%\input sn-article.bbl

\end{document}
